\section{持续学习与多 agent 协作}

\subsection{Multi-Agent 体现与 imitation}

\begin{figure}[H]
\centering
\includegraphics[width=0.8\textwidth]{slide-02-12.jpg}
\caption{Slide 02-12 展示 multi-agent scenario 中的模仿学习与 communication loops 。}
\end{figure}

在 multi-agent 场景下，每个 agent 可以 imitating 不同 specialist，彼此之间通过 gating network share policy logits。CSI gating 允许在 runtime 选择最合适的 expert。

\begin{importantbox}{Multi-agent imitation 的力量}
通过 multi-agent imitation，系统可以在单 agent 模型无法覆盖的复杂任务中分工协作。每个 agent 负责一段 trajectory，最后通过 aggregator 合成一致策略。
\end{importantbox}

\subsection{持续学习与自我改进}

\begin{figure}[H]
\centering
\includegraphics[width=0.8\textwidth]{slide-02-13.jpg}
\caption{Slide 02-13 画出了 continuous improvement loop：data → model → evaluation → deployment。}
\end{figure}

Continuous improvement loop 里，每次 deployment 都需把 human override 作为 new training data，并在 nightly retrain pipeline 中引入 multi-agent critic 作为 auxiliary loss。

\begin{knowledgebox}{Self-supervised critic}
使用 self-supervised critic 可以在缺乏 human comparison 的情况下估计 reward gap，配合 human override 形成 stratified preference signal。
\end{knowledgebox}

\subsection{本章小结}

多 agent imitation 与 continuous improvement loop 让模仿学习变成一个不断自我校准的系统；只要在 evaluation → retrain → deployment 中加上 traction guardrails，就能持续提升。

